\scanpage{48}
\mspage{- IV-977 -22-}

(18.13.4) Le résultat de (18.13.2), \deleted{joint au fait que pour une k-algèbre analytique A, $\overline{\Omega}^1_{A/k}$ soit un A-module de type fini}, peut encore s'exprimer en disant que dans la catégorie des A-modules \uline{de type fini}, $\overline{\Omega}^1_{A/k}$ est un objet \uline{représentant le foncteur covariant} $L\rightsquigarrow\operatorname{Dér}_{\unsure{k|A}}(A,L)$ \inserted{($\uline{0}_{\mathrm{III}}$,8.1.11)}.

\figure{\begin{tikzpicture}\draw (0,0.2) .. controls (-0.1,-0.2) and (0.6,-0.2) .. (1,-0.2)--(4.5,-0.2) .. controls (4.8,-0.2) and (4.8,-0.3) .. (4.8,-0.5);\end{tikzpicture}}

Nous allons généraliser ce résultat. Considérons d'abord un anneau topologique A, et une A-algèbre topologique B qui est un \uline{anneau semi-local noethérien, muni de sa topologie $\underline r$-préadique} ($\uline{r}$ étant le radical de B). Alors, pour tout B-module \uline{de type fini} L, toute A-dérivation D de B dans L est \uline{continue} pour les topologies $\underline r$-préadiques sur B et L, puisque l'on a $D(\underline r^{j+1})\subset\underline r^jL$ pour tout entier j ; en vertu de (\uline{0},20.4.8) elle s'écrit donc $D=u\circ d_{B/A}$, où $u:\Omega^1_{B/A}\to L$ est un homomorphisme de B-modules (nécessairement \uline{continu} pour les topologies $\underline r$-préadiques, et on sait que dans le cas considéré la topologie \inserted{canonique} de $\Omega^1_{B/A}$ est la topologie \deleted{\illegible{2}} \inserted{$r$-préadique} ((\uline{0},20.4.5)). Nous désignerons par $\overline{\Omega}^1_{B/A}$ le quotient de $\Omega^1_{B/A}$ par l'\uline{intersection des noyaux des homomorphismes} $\Omega^1_{B/A}\to L$ pour les B-modules \inserted{L} \uline{de type fini} (ou, ce qui revient au même puisque \deleted{\illegible{1}} L est noethérien, l'intersection des sous-modules M de $\Omega^1_{B/A}$ tels que $\Omega^1_{B/A}/M$ soit \uline{de type fini}). \inserted{On a donc $\operatorname{Hom}_B(\overline{\Omega}^1_{B/A},L)=\operatorname{Hom}_B(\Omega^1_{B/A},L)$ pour tout B-module de type fini L.} Nous désignerons par $\overline d_{B/A}$ ou simplement $\overline d$ l'application composée $B\xrightarrow{d_{B/A}}\Omega^1_{B/A}\xrightarrow{r}\overline{\Omega}^1_{B/A}$ ; c'est une \inserted{A-}dérivation, et $\overline{\Omega}^1_{B/A}$ est engendré par les éléments $\overline d_{B/A}(x)$ où x parcourt B, en vertu de (\uline{0},20.4.7) ; la définition précédente montre que l'application $v\rightsquigarrow v\circ\overline d_{B/A}$ est un isomorphisme de B-modules, fonctoriel en L
\margin{*}
\figure{\begin{tikzpicture}\draw (0.4,0.6) .. controls (-0.2,0.2) and (-0.3,-0.5) .. (0,-0.5) .. controls (1.6,0.1) and (0.7,0.8) .. (0.4,0.6);\draw (0.7,-0.15) .. controls (0.7,-0.7) and (2.4,-0.5) .. (2.4,0.35) .. controls (2.4,0.6) and (2.25,0.7) .. (2.15,0.7);\draw (6,0.7) .. controls (4.7,0.5) and (4.9,-0.3) .. (5.2,-0.3) .. controls (6.6,-0.2) and (6.7,0.6) .. (6,0.7);\end{tikzpicture}}
\margin{à quoi servent ces hypothèses ? Prendre un anneau topologique plus général ?}
\[
(18.13.4.1)\qquad \operatorname{Hom}_B(\overline{\Omega}^1_{B/A},L)\xrightarrow{\sim}\operatorname{Dér}_A(B,L)\inserted{=\operatorname{Dér.cont}_A(B,L)}
\]
L étant toujours supposé \uline{de type fini}. \deleted{\illegible{1}} Lorsque l'on est dans les conditions de (18.13.2), le A-module $\overline{\Omega}^1_{A/k}$ qui vient d'être défini co-
\margin{47}
